\subsection{Solution to Problem 19: Complex Numbers and Equilateral Triangles}

\begin{solution}
\textbf{(i)} Since $w = e^{2\pi i/3}$, we have $w^3 = e^{2\pi i} = 1$, so $w$ is a non-trivial cube root of unity. For $w \ne 1$,
\[
1 + w + w^2 = \frac{1-w^3}{1-w} = \frac{1-1}{1-w} = 0.
\]

\textbf{(ii)} If triangle $ABC$ is anticlockwise and equilateral, a rotation of $120^\circ$ about the centroid $G = \dfrac{a+b+c}{3}$ maps each vertex to the next:
\[
b - G = w(a-G), \qquad c - G = w(b-G).
\]
From the first equation,
\[
b = G + w(a-G) = (1-w)G + wa = \frac{1-w}{3}(a+b+c) + wa.
\]
Rearranging and using $1+w+w^2=0$ (so $1-w = -w-w^2$), we obtain
\[
a + bw + cw^2 = 0.
\]

\textbf{(iii)} For an equilateral triangle, either the anticlockwise condition $a+bw+cw^2=0$ or the clockwise condition $a+bw^2+cw=0$ holds.

Multiplying $a+bw+cw^2=0$ by $w$ gives $aw+bw^2+c=0$, so $c = -aw-bw^2$. Substituting into $a+bw^2+cw=0$:
\[
a + bw^2 + w(-aw-bw^2) = a(1-w^2) + bw^2(1-w) = 0.
\]
Using $1+w+w^2=0$, we have $1-w^2=w$ and $1-w=w^2$, so this becomes $w(a+b)=0$, hence $a+b=-c$ and $a+b+c=0$.

Subtracting $a+b+c=0$ from $a+bw+cw^2=0$:
\[
b(w-1)+c(w^2-1)=0.
\]
Since $w^2-1=(w-1)(w+1)$ and $w+1=-w^2$, we get $(w-1)(b-cw^2)=0$. As $w\ne 1$, it follows that $b=cw^2$; similarly $a=bw^2$ and $c=aw^2$.

From $a+bw+cw^2=0$ and $a+b+c=0$, expand $a^2$, $b^2$, $c^2$ in terms of one another and add. Using $w+w^2=-1$,
\[
a^2+b^2+c^2 = (a^2+b^2+c^2)(w+w^2) + 2(ab+bc+ca)
= -(a^2+b^2+c^2) + 2(ab+bc+ca).
\]
Therefore $2(a^2+b^2+c^2)=2(ab+bc+ca)$, so
\[
a^2+b^2+c^2 = ab+bc+ca.
\]
\end{solution}

\begin{takeaways}
\begin{itemize}
    \item Cube roots of unity: $1+w+w^2=0$ where $w=e^{2\pi i/3}$
    \item Rotation by $120^\circ$ about the centroid maps vertices of an equilateral triangle
    \item Use both orientations to derive $a+b+c=0$, then simplify to the required identity
\end{itemize}
\end{takeaways}
